\section{Related works}
\subsection{Semantic 3D Gaussians Splatting}
Recent advancements have significantly improved semantic understanding in 3D scene representations. Early works build upon implicit radiance field representations such as NeRF. Methods like Semantic NeRF~\cite{semanticnerf}, Panoptic Lifting~\cite{Panopticlifting}, and Contrastive Lift~\cite{ContrastiveLift} pioneered the integration of semantic labels into NeRF-based frameworks, achieving sharp and accurate 3D semantic segmentation. Subsequent approaches—such as Distilled FFD~\cite{FFD}, F3F~\cite{F3F}, panoptic NeRF~\cite{panopticnerf}, NeRF-SOS~\cite{nerfsos}, Interactive Segmentation of Radiance Fields~\cite{ISRF}, Weakly Supervised 3D Open-Vocabulary Segmentation~\cite{weakly}, FFNF~\cite{FFNF}, VL-Fields~\cite{VL-Fields}, Featurenerf~\cite{featurenerf}, Distilled Feature Fields~\cite{DistilledFF}, and LERF~\cite{lerf} leverage semantic features from pretrained vision-language models like CLIP~\cite{CLIP}, LSeg\cite{Lseg}, or DINO\cite{dino} to enable open-vocabulary, pixel-level semantic understanding.

More recently, 3D Gaussian Splatting (3DGS)~\cite{3DGS} has emerged as a compelling explicit alternative to NeRF due to its faster training, real-time rendering capabilities, and high-quality scene reconstruction. Semantic extensions of 3DGS have shown great promise. Semantic Gaussians~\cite{semanticgs} and ConceptFusion~\cite{conceptfusion} project features from 2D pretrained encoders onto 3D Gaussians, enabling open-vocabulary understanding. OpenGaussian~\cite{open} and OpenScene~\cite{openscene} achieve point-level semantic segmentation and produce sharp and accurate semantic segmentation results. Feature 3DGS~\cite{feature3dgs}, CLIP-FO3D~\cite{Clip-fo3d}, and PLA~\cite{pla} further distill semantic features from 2D foundation models into the 3D domain using differentiable rasterization. Thanks to the real-time rendering and optimization capabilities of 3DGS, notable examples like SGS-SLAM~\cite{sgs}, SemGauss-SLAM~\cite{sem-slam}, and GS³LAM~\cite{gs3lam} also integrate semantic features into 3DGS-based SLAM pipelines for robust and scalable scene reconstruction.

Motivated by the strong generalization and ease of training of Feature 3DGS~\cite {feature3dgs}, we adopt its framework as the foundation for our semantic Gaussian Splatting module.

\subsection{Dynamic Gaussians Splatting}
Several recent works have extended 3D scene representation to dynamic scenes by introducing temporal modeling, enabling novel view synthesis in dynamic environments. Based on NeRF, early approaches such as D-NeRF~\cite{DNERF} introduced time as an explicit input to model non-rigid deformations, while Nerfies~\cite{nerfies} leveraged deformation fields with elastic regularization to animate casually captured scenes. HyperNeRF~\cite{HyperNeRF} advanced this by modeling topological changes via a hyperspace embedding, handling complex non-rigid transformations. STaR~\cite{star} decomposes scenes into static and dynamic components using a self-supervised rigid-body model. 

As for 3DGS-based models, 4D Gaussian Splatting (4D-GS)~\cite{4dgs} introduces a unified representation that combines 3D Gaussians with 4D neural voxels with a lightweight MLP to efficiently predict Gaussian deformations over time. Gaussian-Flow~\cite{gs-flow} alternatively models temporal attribute deformations using time and frequency domain parameterizations. Dynamic 3D Gaussians~\cite{dynamic3dgs} uses dense 6-DOF tracking and dynamic reconstruction through enforcing local rigidity constraints to enable Gaussians to maintain persistent attributes when they move and rotate over time. Explicit sampling strategies like Fully Explicit Dynamic Gaussian Splatting~\cite{fully} have also been proposed to separate static and dynamic components during training to represent continuous motion and reduce computational cost.

In our work, we adopt the 4D-GS~\cite{4dgs} framework as the foundation for our dynamic Gaussians module, leveraging its efficient modeling of dynamic scenes and real-time rendering capabilities.



\subsection{Uncertainty and view selection}
Driven by the need to reduce training costs and enable active training, a paradigm where models selectively acquire data based on informativeness, many recent works have focused on extracting useful cues from the trained model, existing data, and candidate viewpoints. A central theme across these methods is the estimation and exploitation of uncertainty. Magic Moments~\cite{MagicMoments} introduces higher moments of rendering equations to quantify the uncertainty of radiance fields and show their relationship with rendered error and their usage as powerful NBV selection criteria. CSS~\cite{CSS} utilizes the expectation and variance of categorical distributions to probabilistically update semantic maps within the 3D Gaussian Splatting (3DGS) framework. Similarly, ~\cite{activeneural} proposes a colorized surface voxel to interpret color uncertainty and surface information. 
While these methods primarily adopt Bayesian formulations to explicitly quantify uncertainty, others infer uncertainty by learning data distributions via alternative models. DNRSelect~\cite{dnrselect} employs a reinforcement learning-based view selector trained on rasterized images to find optimal views for deferred neural rendering. ActiveInitSplat~\cite{activeinitsplat} defines a black-box objective function based on density and occupancy metrics and employs a Gaussian process as the surrogate model to predict informative views.   

Beyond uncertainty estimation, several methods directly target maximizing information gain. GauSS-MI~\cite{gaussmi} utilizes Shannon Mutual Information to select informative viewpoints. FisherRF~\cite{fisherrf} introduces a Fisher Information-based metric to guide view selection in radiance fields without reliance on ground-truth data. AG-SLAM~\cite{agslam} further extends this idea by integrating it into an active SLAM framework, enabling online scene reconstruction and autonomous exploration via information-driven trajectory planning.

In this work, we adopt Fisher Information as the unified selection criterion for both semantic Gaussian Splatting and deformation-based dynamic modeling, facilitating view selection that balances scene semantics and motion-aware updates.





